\begin{ex}
Tờ tiền 100 đồng cộng với tờ tiền 200 đồng bằng bao nhiêu đồng?
\shortans{300}
\loigiai{
Tổng: $100 + 200 = 300$ đồng.
}
\end{ex}

\begin{ex}
Tờ tiền 200 đồng cộng với tờ tiền 500 đồng bằng bao nhiêu đồng?
\shortans{700}
\loigiai{
Tổng: $200 + 500 = 700$ đồng.
}
\end{ex}

\begin{ex}
Có hai tờ tiền 500 đồng. Tổng số tiền là bao nhiêu đồng?
\shortans{1000}
\loigiai{
Tổng: $500 + 500 = 1000$ đồng.
}
\end{ex}

\begin{ex}
Có năm tờ tiền 200 đồng. Tổng số tiền là bao nhiêu đồng?
\shortans{1000}
\loigiai{
Tổng: $200 \times 5 = 1000$ đồng.
}
\end{ex}

\begin{ex}
Có mười tờ tiền 100 đồng. Tổng số tiền là bao nhiêu đồng?
\shortans{1000}
\loigiai{
Tổng: $100 \times 10 = 1000$ đồng.
}
\end{ex}

\begin{ex}
Bạn Mai có 1 tờ 500 đồng và 2 tờ 200 đồng. Mai có tất cả bao nhiêu đồng?
\shortans{900}
\loigiai{
Mai có: $500 + 200 + 200 = 900$ đồng.
}
\end{ex}

\begin{ex}
Bạn Nam có 1 tờ 1000 đồng, Nam mua một cục tẩy hết 600 đồng. Nam còn lại bao nhiêu đồng?
\shortans{400}
\loigiai{
Nam còn: $1000 - 600 = 400$ đồng.
}
\end{ex}

\begin{ex}
Một cây kem giá 800 đồng, bé đưa cho cô bán hàng 1000 đồng. Cô bán hàng phải trả lại bé bao nhiêu đồng?
\shortans{200}
\loigiai{
Trả lại: $1000 - 800 = 200$ đồng.
}
\end{ex}

\begin{ex}
Một phong bì thư giá 300 đồng, một con tem giá 500 đồng. Mua cả phong bì và tem hết bao nhiêu đồng?
\shortans{800}
\loigiai{
Hết: $300 + 500 = 800$ đồng.
}
\end{ex}

\begin{ex}
Lan có 700 đồng, mẹ cho thêm Lan 300 đồng. Lan có tất cả bao nhiêu đồng?
\shortans{1000}
\loigiai{
Lan có: $700 + 300 = 1000$ đồng.
}
\end{ex}

\begin{ex}
Hùng có 900 đồng, Hùng mua một chiếc kẹo mút hết 500 đồng. Hùng còn lại bao nhiêu đồng?
\shortans{400}
\loigiai{
Hùng còn: $900 - 500 = 400$ đồng.
}
\end{ex}

\begin{ex}
Có bao nhiêu tờ tiền 100 đồng để đổi được một tờ tiền 500 đồng?
\shortans{5}
\loigiai{
Cần: $500 : 100 = 5$ tờ.
}
\end{ex}

\begin{ex}
Có bao nhiêu tờ tiền 200 đồng và một tờ 100 đồng để đổi được 500 đồng? Cần mấy tờ 200 đồng?
\shortans{2}
\loigiai{
Vì $200 \times 2 + 100 = 500$ đồng nên cần 2 tờ 200 đồng.
}
\end{ex}

\begin{ex}
Một chiếc bút chì giá 700 đồng, một chiếc thước kẻ giá 300 đồng. Chiếc bút chì đắt hơn chiếc thước kẻ bao nhiêu đồng?
\shortans{400}
\loigiai{
Đắt hơn: $700 - 300 = 400$ đồng.
}
\end{ex}

\begin{ex}
Mua 2 chiếc bánh rán, mỗi chiếc giá 400 đồng. Hết tất cả bao nhiêu đồng?
\shortans{800}
\loigiai{
Hết: $400 \times 2 = 800$ đồng.
}
\end{ex}

\begin{ex}
Một quyển vở nhỏ giá 900 đồng, bạn An chỉ có 600 đồng. An còn thiếu bao nhiêu đồng để mua quyển vở?
\shortans{300}
\loigiai{
Thiếu: $900 - 600 = 300$ đồng.
}
\end{ex}

\begin{ex}
Điền số thích hợp: $400\text{ đồng} + \dots\text{ đồng} = 1000\text{ đồng}$.
\shortans{600}
\loigiai{
Số cần điền: $1000 - 400 = 600$.
}
\end{ex}

\begin{ex}
Điền số thích hợp: $1000\text{ đồng} - \dots\text{ đồng} = 300\text{ đồng}$.
\shortans{700}
\loigiai{
Số cần điền: $1000 - 300 = 700$.
}
\end{ex}

\begin{ex}
Trong ống tiết kiệm của bé có 3 tờ 200 đồng và 4 tờ 100 đồng. Bé đã tiết kiệm được bao nhiêu đồng?
\shortans{1000}
\loigiai{
Tiết kiệm được: $200 \times 3 + 100 \times 4 = 1000$ đồng.
}
\end{ex}

\begin{ex}
Một gói bim bim giá 1000 đồng, bạn Bình có 800 đồng. Bình cần thêm bao nhiêu đồng nữa?
\shortans{200}
\loigiai{
Cần thêm: $1000 - 800 = 200$ đồng.
}
\end{ex}

\begin{ex}
Bác nông dân bán 1 mớ rau muống được 800 đồng, bác mua 1 củ hành hết 300 đồng. Bác còn lại bao nhiêu đồng?
\shortans{500}
\loigiai{
Còn lại: $800 - 300 = 500$ đồng.
}
\end{ex}

\begin{ex}
Lấy số tiền lớn nhất trong các tờ tiền: 100 đồng, 200 đồng, 500 đồng, 1000 đồng. Tờ tiền đó có mệnh giá bao nhiêu đồng?
\shortans{1000}
\loigiai{
Mệnh giá lớn nhất là 1000 đồng.
}
\end{ex}

\begin{ex}
Tờ tiền có mệnh giá nhỏ nhất trong các tờ trên là bao nhiêu đồng?
\shortans{100}
\loigiai{
Mệnh giá nhỏ nhất là 100 đồng.
}
\end{ex}

\begin{ex}
Hiệu giữa tờ tiền 1000 đồng và tờ tiền 500 đồng là bao nhiêu đồng?
\shortans{500}
\loigiai{
Hiệu: $1000 - 500 = 500$ đồng.
}
\end{ex}

\begin{ex}
Tính: $200\text{ đồng} + 300\text{ đồng} + 500\text{ đồng}$.
\shortans{1000}
\loigiai{
Tổng: $200 + 300 + 500 = 1000$ đồng.
}
\end{ex}
